\documentclass[11pt]{scrartcl}

\usepackage{amsmath,amsthm,amssymb,mathtools}
\usepackage{hyperref}
\usepackage{lastpage}
\usepackage{multicol}

\theoremstyle{plain} 
\newtheorem{theorem}{Theorem}[section]
\newtheorem{lemma}[theorem]{Lemma}
\newtheorem*{question}{Question}
\theoremstyle{definition}
\newtheorem*{definition}{Definition} 
\theoremstyle{remark}
\newtheorem*{remark}{Remark}

\usepackage[none]{hyphenat}
\setlength{\parindent}{0pt}
% setlength{\emergencystretch}{3em}

\renewcommand\labelitemi{$\vcenter{\hbox{\tiny$\bullet$}}$} 

\usepackage[margin=1in]{geometry}
\usepackage{fancyhdr}
\pagestyle{fancy}
\lhead{ESC194: Calculus I} 
\rhead{\thepage}
\cfoot{}

\title{ESC194: Calculus I \\[0.1in] Exam Review}
\author{Brendan Tang}
\date{Fall 2026}

\begin{document}

\maketitle
\thispagestyle{empty}

\tableofcontents

\newpage
\section{Limits and Continuity}
\subsection{The Precise Definition of a Limit}

The precise definition of a limit is as follows.

\begin{definition}
    If $\forall\epsilon>0$, $\exists\delta>0$ such that $\lvert x-a \rvert<\delta\implies
    \lvert f(x)-L \rvert<\epsilon$, then \[ \lim_{x \to a}f(x)=L. \]
\end{definition}

For one-sided limits, $x$ is only approaching $a$ from one side. Thus, we define the left-handed
limit as follows.

\begin{definition}
    If $\forall\epsilon>0$, $\exists\delta>0$ such that $a-\delta<x<a\implies \lvert f(x)-L \rvert
    <\epsilon$, then \[ \lim_{ x \to a^- }f(x)=L. \]
\end{definition}

As a number becomes arbitrarily large, it eventually surpasses every real number. Thus, we define
limits to or approaching infinity as follows.

\begin{definition}
    If $\forall M>0$, $\exists\delta>0$ such that $0<\lvert x-a \rvert<\delta\implies f(x)>M$, then
    \[ \lim_{ x \to a }f(x) \to \infty. \]
    Conversely, if $\forall\epsilon>0$, $\exists M>0$ such that $x>M\implies \lvert f(x)-L \rvert
    <\epsilon$, then \[ \lim_{ x \to \infty }f(x)=L. \]
\end{definition}

\begin{question}
    How can we rigorously define right-handed and negative infinite limits? What about limits 
    involving both concepts?
\end{question}

Here are the general steps to a rigorous proof.

\begin{enumerate}
    \item Given $\epsilon>0$, find $\delta>0$. Start with $\lvert f(x)-L \rvert<\epsilon$ and manipulate
    it into the form $\lvert x-a \rvert<\delta_{0}$. This step may require bounding $\delta\leq n$.
    \item \begin{proof} 
        Given $\epsilon$, let $\delta=min(n,\delta_{0})$. Then, show that $\lvert x-a 
        \rvert<\delta\implies \lvert f(x)-L \rvert<\epsilon$, as desired. 
    \end{proof}
\end{enumerate}

\subsection{Limit Theorems}

Here are some important limit theorems.

\begin{theorem}
    If $\lim_{ x \to a }f(x)$ and $\lim_{ x \to a }g(x)$ exist, then
    \begin{itemize}
        \item (Constant LT) $\lim_{ x \to a }c=c$
        \item (Sum LT) $\lim_{ x \to a }[f(x)+g(x)]=\lim_{ x \to a }f(x)+\lim_{ x \to a }g(x)$
        \item (Product LT) $\lim_{ x \to a }[f(x)\cdot g(x)]=\lim_{ x \to a }f(x)\cdot 
        \lim_{ x \to a }g(x)$
        \item (Constant multiple LT) $\lim_{ x \to a }cf(x)=c\lim_{ x \to a }f(x)$
        \item (Power LT) $\lim_{ x \to a }[f(x)]^n=[\lim_{ x \to a }f(x)]^n$
        \item (Root LT) $\lim_{ x \to a }\sqrt[n]{ f(x) }=\sqrt[n]{ \lim_{ x \to a }f(x) }$
    \end{itemize}
\end{theorem}

\begin{theorem}[Direct Substitution]
    If $f$ is continuous at $a$, then \[ \lim_{ x \to a } f(x)=a. \]
\end{theorem}

The following theorem is often used to evaluate trigonometric limits.

\begin{theorem}[Squeeze Theorem]
    If $g(x)<f(x)<h(x)$ and $\lim_{ x \to a }g(x)=\lim_{ x \to a }h(x)=L$, then 
    \[ \lim_{ x \to a } f(x)=L. \]
\end{theorem}

When solving difficult trigonometric limits, it may help to find $\frac{{\sin x}}{x}$ terms.

\subsection{Continuity Theorems}

Here is the definition for continuity.

\begin{definition}
    $f$ is continuous at a point $a$ if $\lim_{ x \to a }f(x)=f(a)$. 
\end{definition}

\begin{question} 
    How can we define continuity from the side? What about continuity on an open or closed interval?
\end{question}

Here are some continuity theorems.

\begin{theorem}
    If $f$ and $g$ are continuous at $a$, then \[ f+g,\ fg,\ cf \] are also continuous at $a$.
\end{theorem}

\begin{theorem}
    The following types of functions are continuous on their domains 
    \begin{multicols}{3}
        \begin{itemize}
            \item polynomial
            \item exponential
            \item rational
            \item logarithmic
            \item root
            \item trigonometric
        \end{itemize}
    \end{multicols}
\end{theorem}

\begin{remark}
    It turns out that all elementary functions are continuous on their domains (why?).
\end{remark}

\newpage
\section{Derivatives}
\subsection{Derivative Theorems}

We begin with some important definitions.

\begin{definition}
    The derivative of $f$ is defined as \[ f'(x)\coloneqq \lim_{ h \to 0 } \frac{{f(x+h)-f(x)}}{h}. \]
    $f$ is differentiable at a point $a$ if $f'(a)$ exists.
\end{definition}

\begin{question} 
    How can we define differentiability from the side? What about differentiability on an open or
    closed interval?
\end{question}

From the definition of the derivative, we see that differentiability at $a$ requires $\lim_{ h \to 0 } 
[f(a+h)-f(a)]=0$ (why?). This implies that $f$ is continuous at $a$. Thus, we arrive at the following theorem.

\begin{theorem}
    If $f$ is differentiable at $a$, then $f$ is continuous at $a$.
\end{theorem}

Be careful, as this theorem does not hold for the converse. That is, continuity at $a$ does not imply
differentiability at $a$. \\[-0.5em]

Here are some useful differentiation theorems.

\begin{theorem}
    If $f'$ exists, then
    \begin{itemize}
        \item (Constant DT) $f(x)=c\implies f'(x)=0$
        \item (Sum DT) $(f+g)'=f'+g'$
        \item (Product DT) $(fg)'=f'g+fg'$
        \item (Reciprocal DT) $\left( \frac{1}{f} \right)'=-\frac{f'}{f^{2}}$
        \item (Quotient DT) $\left( \frac{f}{g} \right)'= \frac{{f'g-fg'}}{g^{2}}$
        \item (Power DT) $f(x)=x^n\implies f'(x)=nx^{n-1}$
        \item (Chain rule) $(f\circ g)'=(f'\circ g)\cdot g'$
    \end{itemize}
\end{theorem}

\subsubsection{Existence Theorems}

These theorems are very important, serving as stepping stones for many proofs that will be encountered in
this course. While their proofs will not be provided here, it is highly encouraged to find and read through 
them yourself.

\begin{theorem}[Intermediate Value Theorem]
    If $f\colon[a,b]\to \mathbb{R}$ is continuous, then for every $N\in\bigl(f(a),f(b)\bigr)$, there exists 
    $c\in(a,b)$ such that $f(c)=N$.
\end{theorem}

\begin{theorem}[Extreme Value Theorem]
    If $f\colon[a,b]\to \mathbb{R}$ is continuous, then there exists $c,d\in[a,b]$ such that $f$ attains an 
    absolute maximum value $f(c)$ and an absolute minimum value $f(d)$.
\end{theorem}

\begin{theorem}[Mean Value Theorem] \label{thm:mvt}
    If $f\colon[a,b]\to \mathbb{R}$ is continuous and differentiable on $(a,b)$, there exists $c\in(a,b)$ 
    such that \[ f'(c)= \frac{{f(b)-f(a)}}{b-a}. \]
\end{theorem}

By using Theorem \ref{thm:mvt} with $f(a)=f(b)$, we get a special case.

\begin{theorem}[Rolle's Theorem] \label{thm:rolle}
    If $f\colon[a,b]\to \mathbb{R}$ is continuous, differentiable on $(a,b)$, and $f(a)=f(b)$, there exists 
    $c\in(a,b)$ such that $f'(c)=0$.
\end{theorem}

\begin{remark}
    Although Theorem \ref{thm:rolle} is presented here as a special case of Theorem \ref{thm:mvt}, 
    chronologically the reverse is true. In fact, Theorem \ref{thm:rolle} is the core tool used to prove 
    Theorem \ref{thm:mvt}, which serves as its generalization.
\end{remark}

\subsubsection{Interlude (Approximation)}

On a bit of an unrelated note, we can use differentials to approximate a function $f$ by the formula
\[ f(x+\Delta x)\approx f(x)+f'(x)\Delta x. \]
This has been included for the sake of completeness, as it has been tested before.

\subsection{Curve Sketching}

By Fermat's theorem on stationary points, we arrive at the following definition.

\begin{definition}
    $\bigl(c,f(c)\bigr)$ is a critical point on $f$ if $f'(c)=0$ or $f'(c)$ does not exist.
\end{definition}

Here are the general steps to curve sketching.

\begin{enumerate}
    \item From $f$:
    \begin{itemize}
        \item Find the domain and range 
        \item Find the $x$ and $y$ intercepts
        \item Find any endpoints or removable discontinuities, if they exist
        \item Find any asymptotes and their behaviour, if they exist
        \item Determine if $f$ has any symmetry
        \item Determine the end behaviour of $f$
    \end{itemize}
    \item From $f'$: 
    \begin{itemize}
        \item Find all critical points on $f$
        \item Determine the intervals of increase/decrease
        \item Find vertical tangents or cusps, if they exist
    \end{itemize}
    \item From $f''$: 
    \begin{itemize}
        \item Find local maxima/minima and any inflection points, if they exist
        \item Determine the intervals of concavity
    \end{itemize}
\end{enumerate}

\subsection{Applications of Derivatives}

One application of derivatives is related rates. Related rates are problems where two or more quantities change,
but are still linked by an equation. To solve a related rates problem, find an equation that directly relates 
the variables and apply the chain rule. \\[-0.5em]

Another application of derivatives is optimization. To solve an optimization problem, reduce the problem to an
equation in one variable, then evaluate at endpoints and critical points.

\newpage
\section{Integrals}
\subsection{Summation}

\end{document}